% =========================================================
% Chapter 2 — Time series decomposition
% Author : Aymen Ben Brik (aymen.benbrik@esprit.tn)
% =========================================================

\chapter{Time series decomposition}
\label{ch:decomposition}

\section*{Introduction}
\addcontentsline{toc}{section}{Introduction}

The first step in time-series analysis is to \emph{decompose} the
observed data into structural components: a slowly varying
\emph{trend} $T_t$, a periodic \emph{seasonality} $S_t$, and a
stationary residual $\varepsilon_t$. The trend and seasonality carry
the deterministic structure; the residual captures the stochastic
fluctuations and is the input to the modelling techniques of
Chapters~\ref{ch:stationarity}--\ref{ch:arch-garch}.

This chapter presents:
\begin{enumerate}
  \item the additive and multiplicative decomposition models;
  \item three trend-estimation strategies (moving average, parametric
        polynomial regression, non-parametric LOESS / splines);
  \item three seasonality-estimation strategies (seasonal averaging,
        dummy regression, cosine--sine Fourier models);
  \item residual diagnostics and the full decomposition workflow.
\end{enumerate}

% =========================================================
\section{The decomposition model}
\label{sec:decomp-model}

\begin{defbox}[Additive vs.\ multiplicative model]
\label{def:decomp}
For an observed time series $(Y_t)_{t=1}^{T}$ with seasonal period $d$:
\[
  \text{Additive:}\quad Y_t = T_t + S_t + \varepsilon_t,
  \qquad
  \text{Multiplicative:}\quad Y_t = T_t \cdot S_t \cdot \varepsilon_t.
\]
Here $T_t$ is a slowly varying \emph{trend}, $S_t$ is a periodic
\emph{seasonal component} of period $d$ (i.e.\ $S_{t + d} = S_t$), and
$\varepsilon_t$ is a zero-mean stationary residual (additive case) or
a positive multiplicative perturbation with $\E[\varepsilon_t] = 1$
(multiplicative case).
\end{defbox}

\begin{rembox}
The two models differ qualitatively:
\begin{itemize}
  \item In the additive model, the seasonal swing has a
        \emph{constant amplitude} regardless of the level of the series.
  \item In the multiplicative model, the swing \emph{grows
        proportionally to $T_t$}.
  \item Taking logarithms converts a multiplicative model into an
        additive one:
        $\log Y_t = \log T_t + \log S_t + \log \varepsilon_t$.
\end{itemize}
This log transform is standard for positive series with growing
amplitude (sales, population, GDP).
\end{rembox}

\paragraph{Identifiability.} Without further constraints the
decomposition is not unique (a constant can move from $T$ to $S$).
Two standard normalisations:
\begin{itemize}
  \item $\sum_{k = 1}^{d} S_k = 0$ in the additive case (seasonality
        sums to zero over a full cycle);
  \item $\sum_{k = 1}^{d} S_k = d$ in the multiplicative case
        (seasonality averages to one).
\end{itemize}

% =========================================================
\section{Trend estimation}
\label{sec:trend}

\subsection{Symmetric (centred) moving average}

\begin{defbox}[Centred moving average]
\label{def:ma}
For an odd window $2k + 1$ with $k \in \N$:
\[
  \widehat{T}_t \;=\; \frac{1}{2k + 1}\sum_{i=-k}^{k} Y_{t+i},
  \qquad t = k+1, \dots, T - k.
\]
Boundary points $t = 1, \dots, k$ and $t = T - k + 1, \dots, T$ are
not estimated.
\end{defbox}

\begin{propbox}[The centred MA preserves linear trends]
\label{prop:ma-linear}
If $Y_t = a + b\,t$ exactly, then $\widehat{T}_t = a + b\,t$ for all
interior $t$.
\end{propbox}

\begin{proof}
Substitute $Y_{t+i} = a + b(t + i)$:
\[
  \widehat{T}_t = \tfrac{1}{2k+1}\sum_{i = -k}^{k}\bigl(a + b(t+i)\bigr)
  = a + b\,t + \tfrac{b}{2k+1}\underbrace{\sum_{i=-k}^{k} i}_{= 0}
  = a + b\,t.
\]
\end{proof}

\begin{propbox}[The centred MA cancels period-$(2k+1)$ seasonality]
\label{prop:ma-season}
If $S_t$ has period $d = 2k + 1$ and $\sum_{k=1}^{d} S_k = 0$, then
$\sum_{i = -k}^{k} S_{t + i} = 0$ for all $t$, hence the MA filters out
$S_t$ exactly.
\end{propbox}

\begin{rembox}
For \emph{even} period $d$ (e.g.\ $d = 12$ for monthly data), one uses
a $2 \times d$ centred MA:
$\widehat{T}_t = \tfrac{1}{2d}\bigl(\tfrac{1}{2}Y_{t-d/2} + Y_{t-d/2+1}
+ \dots + Y_{t+d/2-1} + \tfrac{1}{2}Y_{t+d/2}\bigr)$.
This combines two consecutive symmetric MAs of length $d$ to yield a
window of length $d + 1$ with end-points down-weighted.
\end{rembox}

\begin{exbox}[Numerical illustration]
$Y = (10, 12, 15, 14, 16, 18, 20)$, $k = 1$ (window 3):
\[
  \widehat{T}_2 = \tfrac{37}{3} \approx 12.33, \quad
  \widehat{T}_3 = \tfrac{41}{3} \approx 13.67, \quad
  \widehat{T}_4 = \tfrac{45}{3} = 15.00,
\]
\[
  \widehat{T}_5 = 16.00, \quad
  \widehat{T}_6 = 18.00.
\]
The estimated trend follows the upward drift of the series; $t = 1$
and $t = 7$ are lost.
\end{exbox}

\subsection{Parametric polynomial regression}

\begin{defbox}[Polynomial trend]
$T_t = \beta_0 + \beta_1\, t + \beta_2\, t^2 + \dots + \beta_p\, t^p$
estimated by ordinary least squares (OLS) on $(t, Y_t)$.
\end{defbox}

\begin{rembox}
\begin{itemize}
  \item $p = 1$ : straight line. Use only when the eyeball plot is
        consistent with a linear pattern.
  \item $p = 2, 3$ : flexible enough for most economic / climate data
        on a moderate horizon.
  \item $p \geq 4$ : usually overfit and behave poorly outside the
        sample range. Prefer LOESS or splines.
\end{itemize}
Use AIC / BIC, cross-validation or out-of-sample MSE to select $p$.
\end{rembox}

\subsection{Non-parametric regression}

\begin{defbox}[LOESS / local polynomial regression]
At each $t$ fit a low-degree polynomial on a \emph{neighbourhood} of
$t$, with weights decreasing in $\abs{t' - t}$ (e.g.\ tricube kernel).
The fraction of data used in each local fit is the \emph{span}
parameter $\alpha \in (0, 1)$.
\end{defbox}

\paragraph{Why non-parametric?}
\begin{itemize}
  \item Captures arbitrary smooth trends without imposing a global
        functional form.
  \item Robust variants (\texttt{loess(\dots, family="symmetric")}) are
        less affected by outliers.
  \item Span $\alpha$ small $\Rightarrow$ wiggly fit; large $\Rightarrow$
        over-smooth. Default $\alpha = 0.75$ is often too aggressive
        for time-series trend extraction; $\alpha = 0.1$--$0.3$ is
        usually preferred.
\end{itemize}

\begin{rembox}[Other smoothers]
Cubic smoothing splines minimise
$\sum (Y_t - g(t))^2 + \lambda \int g''(s)^2\,ds$ over twice-differentiable
$g$. The penalty $\lambda$ trades off fidelity vs.\ smoothness;
$\lambda \to 0$ recovers interpolation, $\lambda \to \infty$ a linear
fit. Wavelets offer multi-resolution alternatives but are seldom needed
for a one-step decomposition.
\end{rembox}

% =========================================================
\section{Seasonality estimation}
\label{sec:seasonality}

\subsection{Seasonal averaging}

\begin{defbox}[Seasonal averaging after detrending]
\label{def:seas-avg}
Let $\widetilde{Y}_t = Y_t - \widehat{T}_t$ be the detrended series.
For each season $k = 1, \dots, d$ define
\[
  \widehat{S}_k \;=\; \frac{1}{n_k}\sum_{t : \text{season}(t) = k}
                    \widetilde{Y}_t,
\]
where $n_k$ is the number of observations falling in season $k$.
Re-centre by subtracting the overall mean to ensure $\sum_k \widehat{S}_k = 0$.
\end{defbox}

\begin{exbox}[Monthly seasonal index]
With monthly data ($d = 12$) over 10 years, $n_k = 10$ for each $k$.
The estimate $\widehat{S}_{\text{July}}$ is the average of the
ten July deviations from the trend.
\end{exbox}

\subsection{Dummy-variable regression}

\begin{defbox}[Seasonal dummies]
\label{def:dummies}
Let $D_{k,t} = \mathbf{1}_{\{\text{season}(t) = k\}}$ for $k = 2, \dots, d$.
Estimate by OLS:
\[
  Y_t = \beta_0 + \sum_{k=2}^{d} \beta_k\, D_{k,t} + \varepsilon_t.
\]
The omitted season ($k = 1$) is the \emph{baseline}; $\beta_k$
measures the seasonal deviation of season $k$ from the baseline.
\end{defbox}

\begin{rembox}[Why $d - 1$ dummies?]
Including all $d$ dummies plus an intercept produces perfect
multicollinearity ($\sum_k D_{k,t} = 1$ for all $t$). Drop one season
or drop the intercept.
\end{rembox}

\begin{rembox}[Combining trend and seasonality]
A standard ``one-shot'' specification estimates trend and seasonality
\emph{jointly}:
\[
  Y_t = \beta_0 + \beta_1\, t + \sum_{k=2}^{d} \gamma_k\, D_{k,t} + \varepsilon_t.
\]
This avoids the residual contamination that arises when one estimates
the trend first by MA then the seasonality from the detrended series.
\end{rembox}

\subsection{Cosine--sine (Fourier) model}

\begin{defbox}[Trigonometric seasonality]
\label{def:fourier}
With angular frequency $\omega_d = 2\pi / d$:
\[
  S_t = \sum_{j=1}^{J}\bigl[a_j \cos(j\,\omega_d\, t) + b_j \sin(j\,\omega_d\, t)\bigr],
  \qquad J \leq \lfloor d / 2 \rfloor.
\]
The coefficients $(a_j, b_j)$ are estimated by OLS in the regression
$Y_t = \beta_0 + S_t + \varepsilon_t$ (or jointly with a trend).
\end{defbox}

\begin{propbox}[Equivalence with dummies]
\label{prop:fourier-equiv}
With $J = \lfloor d/2 \rfloor$ harmonics (and the constant), the
trigonometric basis spans the same subspace as the seasonal dummies,
hence yields the same fitted seasonal pattern.
\end{propbox}

\begin{rembox}
The Fourier model is \emph{parsimonious}: with $J = 1$ harmonic only,
seasonal swings are described by 2 parameters instead of $d - 1$. This
helps when $d$ is large (e.g.\ $d = 365$ for daily data with annual
seasonality).
\end{rembox}

% =========================================================
\section{Residuals and diagnostics}
\label{sec:residuals}

After decomposition, the residual is
\[
  \widehat{\varepsilon}_t \;=\; Y_t - \widehat{T}_t - \widehat{S}_t.
\]

\begin{rembox}[What to check]
\begin{itemize}
  \item \textbf{No remaining trend.} Plot $\widehat{\varepsilon}_t$ vs.\ $t$;
        any systematic drift means under-fit trend.
  \item \textbf{No remaining seasonality.} Compute the autocorrelation
        function (Chapter 3) and look for spikes at lag $d$ (and its
        multiples).
  \item \textbf{Constant variance (homoscedasticity).} A funnel-shaped
        plot signals heteroscedasticity, suggesting a multiplicative
        rather than additive model — or ARCH/GARCH effects (Chapter 6).
  \item \textbf{Approximate normality.} Histogram + QQ-plot. Important
        for the validity of $t$-based confidence intervals.
\end{itemize}
\end{rembox}

\begin{warnbox}
A common mistake is to declare the decomposition successful as soon as
the residual ``looks small''. The residual must be \emph{structureless}
— that is, free of autocorrelation, of remaining seasonality and of
heteroscedasticity. The next chapters provide the formal tools to
check this.
\end{warnbox}

% =========================================================
\section{Decomposition workflow}
\label{sec:workflow}

The full pipeline:

\begin{enumerate}
  \item \textbf{Plot} $Y_t$ versus $t$. Decide additive vs.\ multiplicative
        from the visual amplitude pattern.
  \item If multiplicative, apply the log transform.
  \item \textbf{Estimate the trend} $\widehat{T}_t$:
        \begin{itemize}
          \item moving average (centred, $2 \times d$ for even $d$);
          \item polynomial regression ($p = 1, 2$ or $3$);
          \item LOESS / spline.
        \end{itemize}
  \item \textbf{Detrend:} $\widetilde{Y}_t = Y_t - \widehat{T}_t$.
  \item \textbf{Estimate the seasonality} $\widehat{S}_t$ from $\widetilde{Y}_t$:
        \begin{itemize}
          \item seasonal averaging (re-centred);
          \item dummy regression;
          \item cosine--sine model.
        \end{itemize}
  \item \textbf{Residual:}
        $\widehat{\varepsilon}_t = Y_t - \widehat{T}_t - \widehat{S}_t$.
  \item \textbf{Diagnostics} on $\widehat{\varepsilon}_t$ (Section~\ref{sec:residuals}).
  \item If diagnostics pass: model the residual with ARMA (Chapter 5).
\end{enumerate}

% =========================================================
\section*{Summary}
\addcontentsline{toc}{section}{Summary}

\begin{center}
\begin{tabular}{lll}
\toprule
\textbf{Component} & \textbf{Method} & \textbf{Pros / Cons} \\
\midrule
Trend       & Moving average ($2k+1$)      & simple, loses $k$ ends \\
            & Polynomial regression        & global, risk of overfit \\
            & LOESS / splines              & flexible, span tuning \\
\midrule
Seasonality & Seasonal averaging           & simple, no model \\
            & Dummy regression             & joint with trend \\
            & Cosine--sine (Fourier)       & parsimonious, multi-harmonic \\
\bottomrule
\end{tabular}
\end{center}

\medskip
\noindent\textbf{What to remember.}
\begin{itemize}
  \item Decomposition is the \emph{first step} of any time-series
        analysis; the rest of the course works on the residual.
  \item Choose additive or multiplicative based on the amplitude
        pattern of the seasonality.
  \item Estimate trend and seasonality jointly (dummy or Fourier
        regression) when feasible — it controls residual contamination.
\end{itemize}
